% ==========================================================================
% COURS 05 -- PARTIE C : SCHÉMAS-BLOCS, BOUCLAGE ET PERTURBATIONS
% ==========================================================================

\section{Représentation par schéma-blocs}\label{sec:sa-schema-blocs}

Un schéma-blocs représente les relations entre les signaux d'un système. Il permet de conserver une lecture fonctionnelle tout en manipulant les fonctions de transfert.

\SASchemaElementsBase

\begin{itemize}
  \item Un \textbf{bloc} représente une relation $S(p)=H(p)E(p)$.
  \item Un \textbf{comparateur} réalise une somme algébrique.
  \item Un \textbf{point de prélèvement} duplique un signal sans le modifier.
  \item Les flèches imposent le sens de propagation de l'information.
\end{itemize}

\begin{SAattention}
Un schéma-blocs ne représente pas directement la circulation d'énergie. Les signaux peuvent être de natures et d'unités différentes. Chaque bloc doit donc être associé à une entrée, une sortie et une unité cohérentes.
\end{SAattention}

\subsection{Blocs en série, en parallèle et en boucle}

\TLReglesBlocs

Pour des blocs en série :
\[
H_{eq}(p)=\prod_i H_i(p).
\]

Pour des blocs en parallèle soumis à la même entrée :
\[
H_{eq}(p)=H_1(p)+H_2(p)
\]
avec les signes indiqués au sommateur.

Pour une rétroaction négative de chaîne directe $D(p)$ et de retour $R(p)$ :
\[
\boxed{H_{BF}(p)=\frac{D(p)}{1+D(p)R(p)}}.
\]
Pour une rétroaction positive, le signe du dénominateur devient négatif.

\subsection{Déplacement des sommateurs et points de prélèvement}

Lorsqu'un sommateur ou un point de prélèvement est déplacé de part et d'autre d'un bloc, il faut introduire le bloc ou son inverse afin de conserver les mêmes relations entre signaux.

\begin{SAmethode}[title={Simplifier un schéma-blocs}]
\begin{enumerate}
  \item Identifier les boucles internes et les réduire en premier.
  \item Regrouper les blocs en série ou en parallèle.
  \item Déplacer un sommateur ou un prélèvement uniquement si cela simplifie réellement la structure.
  \item Vérifier le résultat sur les dimensions et sur des cas limites simples.
\end{enumerate}
\end{SAmethode}

\section{Fonctions de transfert en boucle ouverte et fermée}\label{sec:sa-ftbo-ftbf}

Dans une boucle à retour unitaire :
\[
\varepsilon(p)=E(p)-S(p),\qquad S(p)=D(p)\varepsilon(p).
\]
On obtient :
\[
\frac{S(p)}{E(p)}=\frac{D(p)}{1+D(p)},
\qquad
\frac{\varepsilon(p)}{E(p)}=\frac{1}{1+D(p)}.
\]

Dans le cas général, la fonction de transfert en boucle ouverte est :
\[
L(p)=D(p)R(p).
\]
La fonction de transfert en boucle fermée est :
\[
H_{BF}(p)=\frac{D(p)}{1+L(p)}.
\]

Le polynôme $1+L(p)$ est l'équation caractéristique de la boucle fermée. Ses racines sont les pôles du système bouclé.

\begin{SAapplication}[title={Capteur non unitaire}]
Si la chaîne directe vaut $C(p)P(p)$ et le capteur $H(p)$ :
\[
\frac{S(p)}{E(p)}=\frac{C(p)P(p)}{1+C(p)P(p)H(p)}.
\]
La sortie ne suit pas nécessairement la consigne avec un gain statique égal à 1. L'étalonnage du capteur et les unités du comparateur doivent être pris en compte.
\end{SAapplication}

\section{Systèmes perturbés}\label{sec:sa-perturbations}

\TLSchemaPerturbe

Une perturbation peut intervenir à différents endroits de la chaîne. Son effet sur la sortie dépend de son point d'application.

\subsection{Théorème de superposition}

Pour un système linéaire à plusieurs entrées, la sortie totale est la somme des sorties dues à chaque entrée prise séparément :
\[
S(p)=H_E(p)E(p)+H_Q(p)Q(p).
\]

Pour déterminer $H_E$, on annule les autres entrées. Pour déterminer $H_Q$, on annule la consigne.

\begin{SAapplication}[title={Perturbation appliquée à l'entrée du procédé}]
Avec une chaîne directe $C(p)P(p)$, un retour unitaire et une perturbation $Q(p)$ ajoutée avant $P(p)$ :
\[
S(p)=\frac{C(p)P(p)}{1+C(p)P(p)}E(p)
+\frac{P(p)}{1+C(p)P(p)}Q(p).
\]
Un grand gain de boucle réduit l'influence de la perturbation dans la bande où $|C(j\omega)P(j\omega)|$ est grand.
\end{SAapplication}

\subsection{Sensibilité et rejet de perturbation}

On définit la fonction de sensibilité :
\[
S_e(p)=\frac{1}{1+L(p)},
\]
et la sensibilité complémentaire :
\[
T(p)=\frac{L(p)}{1+L(p)}.
\]
Elles vérifient :
\[
S_e(p)+T(p)=1.
\]

À basse fréquence, un grand gain de boucle rend $S_e$ faible : l'erreur et certaines perturbations sont rejetées. À haute fréquence, un gain trop grand amplifie le bruit de mesure et peut dégrader la robustesse.

\section{Précision statique et classe d'un système}\label{sec:sa-precision-statique}

Pour un retour unitaire :
\[
\varepsilon(p)=\frac{E(p)}{1+L(p)}.
\]
L'erreur statique est obtenue par le théorème de la valeur finale :
\[
\varepsilon_s=\lim_{p\to0}p\varepsilon(p).
\]

On écrit la boucle ouverte sous la forme :
\[
L(p)=\frac{K}{p^\alpha}\,L_0(p),\qquad L_0(0)=1.
\]
L'entier $\alpha$ est la \textbf{classe} du système.

\subsection{Erreur pour une consigne échelon}

Pour $E(p)=E_0/p$ :
\[
\varepsilon_s=\lim_{p\to0}\frac{E_0}{1+L(p)}.
\]
Ainsi :
\begin{itemize}
  \item classe 0 : $\varepsilon_s=E_0/(1+K)$ si $L(0)=K$ ;
  \item classe $\geq1$ : $\varepsilon_s=0$ si le système bouclé est stable.
\end{itemize}

\subsection{Erreur pour une consigne rampe}

Pour $E(p)=V_0/p^2$ :
\[
\varepsilon_v=\lim_{p\to0}\frac{V_0}{p(1+L(p))}.
\]
Ainsi :
\begin{itemize}
  \item classe 0 : erreur infinie ;
  \item classe 1 : erreur finie $V_0/K_v$ ;
  \item classe $\geq2$ : erreur nulle.
\end{itemize}

\begin{SAattention}
Ajouter un intégrateur améliore la précision à basse fréquence mais diminue généralement les marges de stabilité. La précision ne peut donc pas être améliorée indépendamment de la stabilité et de la rapidité.
\end{SAattention}

\section{Identifier un modèle de comportement}\label{sec:sa-identification}

La modélisation de comportement s'appuie sur des essais. Elle est utile lorsque les paramètres physiques sont inconnus ou lorsque le modèle de connaissance serait trop complexe.

\begin{SAmethode}[title={Démarche d'identification temporelle}]
\begin{enumerate}
  \item Placer le système dans un régime initial stable.
  \item Appliquer une entrée connue, souvent un échelon.
  \item Enregistrer l'entrée et la sortie avec une fréquence d'échantillonnage suffisante.
  \item Retirer les offsets et repérer un éventuel retard.
  \item Choisir une famille de modèle : premier ordre, second ordre, intégrateur, retard.
  \item Estimer les paramètres puis comparer la réponse simulée aux mesures.
\end{enumerate}
\end{SAmethode}

\subsection{Identification d'un premier ordre}

Pour un modèle $K/(1+\tau p)$ :
\[
K=\frac{\Delta s_\infty}{\Delta e},
\]
et $\tau$ est lu à $63{,}2\%$ de la variation finale. En présence d'un retard $T$, on utilise :
\[
H(p)=\frac{K e^{-Tp}}{1+\tau p}.
\]

\subsection{Identification d'un second ordre pseudo-périodique}

Pour une réponse présentant des dépassements :
\[
D_1=\exp\left(-\frac{\pi\xi}{\sqrt{1-\xi^2}}\right),
\]
d'où :
\[
\xi=\frac{-\ln D_1}{\sqrt{\pi^2+(\ln D_1)^2}}.
\]
La pseudo-période mesurée $T_a$ donne :
\[
\omega_a=\frac{2\pi}{T_a},\qquad
\omega_0=\frac{\omega_a}{\sqrt{1-\xi^2}}.
\]

\subsection{Validation}

Un modèle identifié doit être validé sur un essai différent de celui utilisé pour l'estimation. On compare notamment :
\begin{itemize}
  \item la valeur finale ;
  \item le temps de réponse ;
  \item les dépassements et la pseudo-période ;
  \item l'erreur quadratique entre mesure et simulation ;
  \item la cohérence des paramètres avec la physique du système.
\end{itemize}
